\chapter{Quantum Error Correction}

Real devices decohere. Error correction protects \emph{logical} qubits using redundancy in \emph{physical} qubits.

\section{Noise and Decoherence}

Errors include bit flips ($X$), phase flips ($Z$), depolarizing noise, relaxation ($T_1$), and dephasing ($T_2$). As of the \cutoff{} research cutoff, large-scale fault-tolerant quantum computation remains an active engineering goal; recent experiments demonstrate improving logical error rates with surface-code scaling \cite{acharya2024}, while commercial roadmaps should be distinguished from established theorems.

\section{Three-Qubit Bit-Flip Code}

Logical encoding: $\ket{0_L}=\ket{000}$, $\ket{1_L}=\ket{111}$. An arbitrary logical qubit $\alpha\ket{0_L}+\beta\ket{1_L}$ can suffer a single $X$ error on one wire; parity measurements reveal the syndrome without collapsing $\alpha,\beta$.

\section{Phase-Flip Code}

Phase errors flip $\ket{+}\leftrightarrow\ket{-}$. Conjugation by $H$ exchanges $X$ and $Z$, reducing phase-flip correction to bit-flip correction in the rotated basis.

\section{Shor Nine-Qubit Code}

The Shor code combines bit- and phase-flip protection using nine physical qubits per logical qubit. Modern research emphasizes surface codes and other topological schemes with higher thresholds and more hardware-friendly layouts.

\section{Error Mitigation vs.\ Correction}

In the NISQ era \cite{preskill2018}, error \emph{mitigation} techniques (zero-noise extrapolation, probabilistic error cancellation) reduce bias without full logical encoding. Mitigation is not a substitute for fault tolerance at scale.

\begin{center}
\fbox{\parbox{0.9\textwidth}{\textbf{Interactive Lab 14:} Inject bit-flip errors and correct with syndromes.}}
\end{center}

\section*{Exercises}
\begin{enumerate}[label=\arabic*.]
\item What syndrome indicates an $X$ error on qubit 2 in the three-bit code?
\item Why cannot you clone $\ket{\psi}$ to protect it?
\item Distinguish physical and logical qubits in one paragraph.
\end{enumerate}
